\section{Introduction}
\label{sec:intro}

A deployed language model may be asked to remove a person's private records,
a copyrighted document, or a harmful capability long after training has
finished. Retraining from scratch for every such request is usually
impractical, so machine unlearning suppresses the designated influence
through targeted post-training updates while preserving the model's
remaining behavior \cite{cao2015towards,bourtoule2021machine,
liu2024rethinking}. The selective requirement is what makes unlearning
hard: forget enough to satisfy the request, but not so aggressively that
useful knowledge is destroyed with it.

The resulting collateral damage is easy to miss because it concentrates.
A model can preserve most retained examples while failing sharply on a
small, coherent subset -- other authors when one author is deleted, benign
science that shares terminology with a suppressed hazardous procedure.
Split-aware evaluations expose this by reporting forget-adjacent retention
separately \cite{cheng2026entanglement,chang2025retain}, but they are
retrospective: the vulnerable behavior is identified only after unlearning
has been run and paid for. Practitioners would rather know \emph{before}
unlearning where damage will land, so that a limited protection or repair
budget can be spent where it is needed.

Prior attempts at such prospective prediction share a hidden assumption:
that susceptibility is a property of the model and data alone, so that some
single score -- lexical or embedding similarity to the forget set, gradient
alignment, gradient magnitude -- should rank vulnerable candidates
regardless of which unlearning method is applied. Empirically these scores
behave inconsistently across methods, and each new proposal appears to win
on some objectives and lose on others. This paper explains that
inconsistency and turns it into structure:

\noindent\emph{Retain-side susceptibility is not a model-intrinsic scalar.
It is conditional on the damage channel induced by the unlearning objective
-- declared in advance by where its loss reads the model.}

\begin{figure*}[t]
\centering
% Figure 1: objective-conditioned damage -> sealed profiles -> crossed test.
% Candidate intensities are deliberately schematic.  The palette is
% grayscale plus one muted blue-gray accent so the logic survives printing.
\definecolor{accent}{HTML}{5E7180}
\definecolor{ink}{HTML}{333333}

\begin{tikzpicture}[
  x=1cm,y=1cm,
  font=\sffamily\scriptsize,
  >={Latex[length=1.65mm,width=1.05mm]},
  paneltitle/.style={font=\sffamily\scriptsize\bfseries,text=ink,anchor=west},
  panelsub/.style={font=\sffamily\scriptsize,text=gray!78!black,anchor=west},
  objective/.style={draw=gray!54,rounded corners=.65mm,fill=white,
                    minimum width=1.58cm,minimum height=1.02cm,line width=.52pt},
  profile/.style={draw=gray!54,rounded corners=.65mm,fill=white,
                  minimum width=1.42cm,minimum height=.98cm,line width=.52pt,
                  align=center},
  candidate/.style={draw=gray!48,rounded corners=.35mm,fill=white,
                    minimum width=.44cm,minimum height=.60cm,inner sep=0pt,
                    line width=.42pt,font=\sffamily\tiny,text=ink},
  seal/.style={draw=accent,rounded corners=.35mm,fill=accent!5,
               inner xsep=.65mm,inner ysep=.45mm,line width=.52pt,
               font=\sffamily\tiny\bfseries,text=accent},
  flow/.style={->,draw=gray!70!black,line width=.62pt},
  trajectory/.style={->,draw=ink!82,line width=.64pt},
  cell/.style={draw=gray!48,fill=gray!3,minimum width=1.00cm,
               minimum height=.63cm,inner sep=0pt,line width=.48pt}
]

\path[use as bounding box] (0,.16) rectangle (17.35,4.72);

% Quiet panel separators.
\draw[gray!25,line width=.4pt] (5.55,.34) -- (5.55,4.58);
\draw[gray!25,line width=.4pt] (11.15,.34) -- (11.15,4.58);

% -------------------------------------------------------------------------
% (a) The same candidates can face different objective-conditioned damage.
\node[paneltitle] at (.16,4.50) {(a)\quad OBJECTIVE-CONDITIONED DAMAGE};
\node[panelsub] at (.16,4.10)
  {same $\theta_0$ and $\mathcal C$; different damage futures};

% Output-reading objective.
\node[objective] (uout) at (.98,3.14) {};
\node[font=\sffamily\tiny\bfseries,text=ink] at (.98,3.48)
  {$\mathcal U_{\mathrm{out}}$};
\foreach \x/\h in {.55/.12,.70/.22,.85/.31,1.00/.18}
  \draw[accent,line width=1.15pt] (\x,2.78) -- (\x,{2.78+\h});
\draw[gray!44,line width=.35pt] (.48,2.78) -- (1.07,2.78);
\node[anchor=west,align=left,font=\sffamily\tiny,text=ink] at (1.12,2.97)
  {loss reads\\[-.35mm]\textbf{outputs}};

% Representation-reading objective.
\node[objective] (urep) at (.98,1.54) {};
\node[font=\sffamily\tiny\bfseries,text=ink] at (.98,1.88)
  {$\mathcal U_{\mathrm{rep}}$};
\foreach \x in {.51,.68,.85}
  \draw[draw=accent,fill=white,rounded corners=.12mm,line width=.4pt]
    (\x,1.18) rectangle ++(.11,.39);
\foreach \y in {1.25,1.37,1.49}
  \fill[accent] (1.02,\y) circle (.75pt);
\node[anchor=west,align=left,font=\sffamily\tiny,text=ink] at (1.13,1.37)
  {loss reads\\[-.35mm]\textbf{hidden states}};

\draw[flow] (uout.east) -- (2.05,3.14);
\draw[flow] (urep.east) -- (2.05,1.54);

% Identical candidate locations; mini-bar height changes the qualitative
% ordering.  The bars carry no numeric experimental value.
\foreach \x/\lab/\height in {2.30/1/.20,2.86/2/.06,3.42/3/.15,3.98/4/.04,4.54/5/.10}{
  \node[candidate] at (\x,3.14) {$x_{\lab}$};
  \fill[gray!68] ({\x-.065},2.85) rectangle ({\x+.065},{2.85+\height});
}
\foreach \x/\lab/\height in {2.30/1/.06,2.86/2/.20,3.42/3/.04,3.98/4/.15,4.54/5/.10}{
  \node[candidate] at (\x,1.54) {$x_{\lab}$};
  \fill[gray!68] ({\x-.065},1.25) rectangle ({\x+.065},{1.25+\height});
}
\node[font=\sffamily\tiny\bfseries,text=ink] at (3.36,.66)
  {same candidates, different damage orderings};

% -------------------------------------------------------------------------
% (b) Both profiles are prospective and use the same candidate universe.
\node[paneltitle] at (5.78,4.50) {(b)\quad PROFILES BEFORE UNLEARNING};
\node[panelsub] at (5.78,4.10)
  {same $\mathcal C$; computed at $\theta_0$ and sealed};

\node[profile] (qg) at (6.55,3.14) {};
\node[font=\sffamily\tiny\bfseries,text=ink] at (6.55,3.48)
  {$\widehat q_G$};
\node[font=\sffamily\tiny,text=ink] at (6.55,3.15)
  {$\ell^-\!\leftarrow\!\theta_0\!\rightarrow\!\ell^+$};
\node[font=\sffamily\tiny,text=ink] at (6.55,2.84)
  {shake energy};

\node[profile] (qh) at (6.55,1.54) {};
\node[font=\sffamily\tiny\bfseries,text=ink] at (6.55,1.88)
  {$\widehat q_H$};
\node[circle,draw=accent,fill=white,minimum size=3.0mm,inner sep=0pt,
      line width=.42pt,font=\sffamily\tiny,text=ink] at (6.55,1.56) {$x$};
\foreach \x/\y in {6.27/1.66,6.82/1.68,6.79/1.39}
  \node[rectangle,draw=accent,fill=accent,minimum size=2.6pt,inner sep=0pt]
    at (\x,\y) {};
\node[font=\sffamily\tiny,text=ink] at (6.55,1.22)
  {hidden neighbors};

\node[font=\sffamily\scriptsize,text=gray!78!black] at (8.82,3.72)
  {shared $\pm$ loss shakes; no candidate backward};
\node[font=\sffamily\scriptsize,text=gray!78!black] at (8.82,2.12)
  {nearest forget-state neighbors};
\draw[flow] (qg.east) -- (7.35,3.14);
\draw[flow] (qh.east) -- (7.35,1.54);

% Accent mini-bars denote prospective exposure, not observed damage.
\foreach \x/\lab/\height in {7.58/1/.18,8.11/2/.07,8.64/3/.15,9.17/4/.05,9.70/5/.10}{
  \node[candidate] at (\x,3.14) {$x_{\lab}$};
  \fill[accent!82] ({\x-.065},2.85) rectangle ({\x+.065},{2.85+\height});
}
\foreach \x/\lab/\height in {7.58/1/.07,8.11/2/.18,8.64/3/.05,9.17/4/.15,9.70/5/.09}{
  \node[candidate] at (\x,1.54) {$x_{\lab}$};
  \fill[accent!82] ({\x-.065},1.25) rectangle ({\x+.065},{1.25+\height});
}
\node[seal] at (10.48,3.14) {SEALED};
\node[seal] at (10.48,1.54) {SEALED};
\node[font=\sffamily\tiny\bfseries,text=ink] at (8.48,.66)
  {two prospective hypotheses, no outcome access};

% -------------------------------------------------------------------------
% (c) Independent trajectories populate a crossed, preregistered endpoint.
\node[paneltitle] at (11.38,4.50) {(c)\quad PREREGISTERED CROSSED TEST};
\node[panelsub] at (11.38,4.10)
  {test, not result: does the relative winner flip?};

% The target trajectories never consume the prospective profiles.
\node[draw=gray!58,rounded corners=.45mm,fill=white,minimum width=.62cm,
      minimum height=.48cm,line width=.5pt,font=\sffamily\tiny\bfseries]
      (theta) at (11.82,2.37) {$\theta_0$};
\node[font=\sffamily\scriptsize,text=gray!78!black] at (12.65,3.72)
  {independent runs};
\node[draw=gray!55,rounded corners=.4mm,fill=white,minimum width=.62cm,
      minimum height=.46cm,line width=.48pt,font=\sffamily\tiny]
      (co) at (12.62,3.12) {$\mathcal U_{\mathrm{out}}$};
\node[draw=gray!55,rounded corners=.4mm,fill=white,minimum width=.62cm,
      minimum height=.46cm,line width=.48pt,font=\sffamily\tiny]
      (cr) at (12.62,1.62) {$\mathcal U_{\mathrm{rep}}$};
\node[draw=gray!55,rounded corners=.4mm,fill=gray!4,minimum width=.58cm,
      minimum height=.46cm,line width=.48pt,font=\sffamily\tiny]
      (do) at (13.57,3.12) {$d_{\mathrm{out}}$};
\node[draw=gray!55,rounded corners=.4mm,fill=gray!4,minimum width=.58cm,
      minimum height=.46cm,line width=.48pt,font=\sffamily\tiny]
      (dr) at (13.57,1.62) {$d_{\mathrm{rep}}$};
\draw[trajectory] (theta.east) -- (12.14,2.37) |- (co.west);
\draw[trajectory] (theta.east) -- (12.14,2.37) |- (cr.west);
\draw[trajectory] (co) -- (do);
\draw[trajectory] (cr) -- (dr);

% Correlation matrix.  All four cells remain visually neutral because the
% preregistered interaction is a question, not an observed result.
\node[font=\sffamily\scriptsize,text=gray!78!black,align=center] at (15.23,3.57)
  {output\\damage};
\node[font=\sffamily\scriptsize,text=gray!78!black,align=center] at (16.39,3.57)
  {representation\\damage};
\node[font=\sffamily\tiny\bfseries,text=ink,align=right] at (14.42,2.82)
  {$\widehat q_G$};
\node[font=\sffamily\tiny\bfseries,text=ink,align=right] at (14.42,2.02)
  {$\widehat q_H$};
\node[cell]        at (15.23,2.82) {$\rho_{G,\mathrm{out}}$};
\node[cell]        at (16.39,2.82) {$\rho_{G,\mathrm{rep}}$};
\node[cell]        at (15.23,2.02) {$\rho_{H,\mathrm{out}}$};
\node[cell]        at (16.39,2.02) {$\rho_{H,\mathrm{rep}}$};

\node[font=\sffamily\scriptsize,text=ink] at (15.17,.79)
  {$\displaystyle
    \Delta=
    (\rho_{G,\mathrm{out}}-\rho_{H,\mathrm{out}})
    -(\rho_{G,\mathrm{rep}}-\rho_{H,\mathrm{rep}})>0\;?$};

\end{tikzpicture}

\caption{\textbf{Susceptibility is tested as an objective-conditioned
interaction.} \textbf{(a)} For the objectives studied, the direct
forget-suppression term reads either outputs or hidden states; the same
$\theta_0$ and retained candidates may therefore yield different damage
rankings. \textbf{(b)} The loss-shake gradient-energy and hidden-state-proximity
hypotheses score the same candidates at $\theta_0$ and are sealed before any
outcome is observed. \textbf{(c)} Independent unlearning trajectories that
never access the scores produce the two damage columns; matrix entries are
Spearman correlations. The preregistered endpoint asks whether the relative
predictive advantage crosses over between channels, $\Delta>0$, not whether
one probe is universally superior. Candidate intensities are schematic, not
experimental results.}
\Description{A three-panel schematic with no observed results. Panel (a)
applies an output-reading and a hidden-state-reading unlearning objective to
the same initial model and retained-candidate set; qualitative bands under
the repeated candidate identifiers illustrate two possible damage orderings.
Panel (b) applies loss-shake gradient-energy and hidden-state-proximity
profiles to that same candidate set and marks both profiles as sealed before
unlearning. Panel (c) shows two independent objective trajectories producing
output-channel and representation-channel damage, followed by a two-by-two
matrix of correlations between each sealed profile and each damage outcome.
The difference-in-differences expression below the matrix asks whether the
relative predictive winner crosses over between channels.}
\label{fig:channel-mechanism}
\end{figure*}

To first order, the damage to a retained candidate $x$ under a parameter
update $\Delta\theta$ is $d(x)\approx g_x^{\top}\Delta\theta$, which makes
$\lVert g_x\rVert$ a direction-free exposure envelope rather than a
guarantee of realized damage. What
$\Delta\theta$ looks like is determined by the objective's loss. Objectives
whose loss is a function of model \emph{outputs} -- token likelihoods or
preferences, as in GA, GradDiff, NPO, SimNPO, IdkDPO, and GRU
\cite{yao2024large,zhang2024negative,fan2024simplicity,maini2024tofu,
wang2025gru} -- motivate gradient-energy exposure. Objectives whose loss is a
function of internal \emph{representations} -- hidden states, as in RMU,
RepNoise, and Circuit Breakers
\cite{li2024wmdp,rosati2024repnoise,zou2024circuit} -- displace activation
geometry and motivate proximity to forget representations. The two
channels therefore call for different predictor families. We test that
claim at the family level with a frozen objective-by-probe interaction;
individual objectives may show partial, mixed, or reversed signal.

Acting on this structure is only useful if the matching signal is cheap at
candidate-pool scale. Exact per-candidate gradient norms require one
backward pass per candidate. We instead estimate the entire pool's gradient
magnitudes with a \emph{backward-free randomized-sensitivity probe}:
symmetric forward evaluations at $\theta_0\pm\eta v_r$ along a few shared
random directions $v_r$ turn squared central differences into an unbiased
(up to $O(\eta^2)$) estimate of every candidate's gradient energy
simultaneously, with Monte Carlo variance controlled by the number of
shared directions. A
fidelity gate separates the estimator's numerical validity from its
predictive value. The representation channel uses a complementary
forward-only hidden-state proximity profile.

Prediction becomes protection through routing. Given a deletion request and
a declared objective, the channel determines the probe; the probe ranks the
discovery-fold candidates; the top-ranked pool receives a fixed repair
budget under a guarded procedure that caps forgetting regression. Our
success criterion is deliberately not a leaderboard: it is the
\emph{interaction} in a preregistered crossed design. For an output-channel
parent (GradDiff) and a representation-channel parent (RMU), we cross each
with matched, mismatched, random, and no-protection selectors and test
whether the channel-matched selector yields the lowest audit-tail damage at
matched forgetting -- per parent. All arms and reach failures are reported;
the channel claim is not declared successful unless both parent-wise tests
pass.

Our contributions are threefold:
\begin{itemize}
    \item \textbf{Channel-conditioned susceptibility.} We test whether the
    predictive signal changes with the objective's channel, declared from
    its loss definition. Evidence comes from a frozen family-level
    interaction over a channel-stratified objective roster, not a
    per-method anecdote.

    \item \textbf{A backward-free probe for the gradient channel.} A
    randomized symmetric-perturbation estimator profiles gradient
    magnitude over the whole candidate pool using only $2R$ batched
    forward sweeps, with a preregistered fidelity gate that certifies when
    the estimate is trustworthy under a given radius and precision.

    \item \textbf{A falsifiable protection test.} A crossed design routes
    the same repair budget through matched, mismatched, random, and no-score
    selectors, testing whether the pre-unlearning channel distinction is
    operationally useful at matched forgetting.
\end{itemize}

We position this work as a prediction-and-protection \emph{layer} on top of
existing unlearners, not as a competing unlearning objective. The remainder
of the paper follows this structure. Section~\ref{sec:related} organizes
prior unlearning methods by channel and identifies the channel-matching
gap. Section~\ref{sec:setup} formalizes damage channels, the prediction
and protection endpoints, and the preregistered interaction criterion
(Figure~\ref{fig:channel-mechanism}). Section~\ref{sec:method} presents
the probes, the fidelity gate, and the guarded channel-matched protection
procedure. Section~\ref{sec:experiments} reports the channel matrix, the
crossed protection experiment, cost, and cross-dataset replication.
